\documentclass[14pt,a4paper]{extarticle}

% Поля страницы
\usepackage[
    left=30mm,
    right=15mm,
    top=20mm,
    bottom=20mm
]{geometry}

\usepackage{fontspec}
\usepackage{polyglossia}
\setmainlanguage{russian}
\setotherlanguage{english}
\setmainfont{Times New Roman}

\usepackage{setspace}
\onehalfspacing

\usepackage{indentfirst}
\setlength{\parindent}{1.25cm}

\usepackage{amsmath, amssymb}
\usepackage{graphicx}
\usepackage{booktabs}
\usepackage{longtable}

\usepackage{hyperref}
\hypersetup{
    colorlinks=true,
    linkcolor=black,
    citecolor=black,
    urlcolor=black
}

\usepackage[
    backend=biber,
    style=gost-numeric,
    language=auto,
    autolang=other
]{biblatex}
\addbibresource{refs.bib}

\begin{document}

\begin{titlepage}
    \centering

    {\large Название университета / института\par}
    \vspace{2cm}

    {\Large Курсовая работа\par}
    \vspace{1cm}

    {\LARGE\bfseries Построение системы многошагового поиска обзорного материала по новостной статье\par}

    \vfill

    \begin{flushright}
        Выполнил: ФИО \\
        Научный руководитель: ФИО
    \end{flushright}

    \vspace{2cm}

    Москва / Город \\
    2026

\end{titlepage}

\tableofcontents
\newpage

\section{Введение}

В условиях быстрого роста объёмов текстовой информации задача поиска и анализа релевантных материалов становится всё более значимой. 
Пользователь, сталкивающийся с новостной статьёй, аналитической заметкой или другим информационным сообщением, часто заинтересован в более широком контексте, связанным с этой темой:
с какими событиями и процессами связаны темы этой статьи, какие выводы можно сделать на основе найденных дополнительных материалов. 
Обычный поиск по ключевым словам в такой ситуации не всегда оказывается достаточным, 
поскольку такой поиск находит лишь близкие по словам и семантике документы, в то время как действительно полезные документы могут быть
связаны, например, логически и не содержать пересекающихся тем.

Особенно важной становится задача многошагового поиска, при котором процесс получения информации не ограничивается одним запросом. 
Вместо этого система последовательно выделяет ключевые темы из исходного текста, расширяет их за счёт связанных понятий, событий и контекстов, 
а затем использует найденную информацию для построения более полного представления о рассматриваемой проблеме. 
Такой подход позволяет переходить от поверхностного поиска по совпадению слов к поиску, ориентированному на смысловые связи между темами.
В отличие от классического информационного поиска, где основное внимание уделяется ранжированию документов по релевантности запросу, 
в многошаговом сценарии существенную роль играет промежуточное представление тем, их развитие и интерпретация найденных связей.

В рамках данной работы рассматривается следующая задача.
Предполагается, что на вход системе подаётся текст (новость/тема/факт из википедии), после чего из него выделяются основные тематики. 
Для каждой тематики выполняется поиск связанных направлений, формируется набор дополнительных материалов, 
а затем строится описание связи между исходной статьёй и найденными темами. 
Такой процесс позволяет получить обзорный реферат, описывающий выводы по главному запросу и связанным подзапросам.


\section{}
\newpage
\printbibliography

\end{document}